\subsection{Definition of symmetric tensors}

Let M be an A-module. We recall (III, p.~501) that \(\mathfrak{S}_n\) operates on the left of the A-module \(\mathbf{T}^n (\mathbf{M})\) , in such a way that

\[
\sigma \left(x _ {1} \otimes x _ {2} \otimes \dots \otimes x _ {n}\right) = x _ {\sigma^ {- 1} (1)} \otimes x _ {\sigma^ {- 1} (2)} \otimes \dots \otimes x _ {\sigma^ {- 1} (n)}
\]

for any \(x_{1},\ldots,x_{n}\in\mathbf{M}\) and \(\sigma\in\mathfrak{S}_{n}\) . The elements \(z\in\mathbf{T}^{n}(\mathbf{M})\) such that \(\sigma\cdot z=z\) for all \(\sigma\in\mathfrak{S}_{n}\) are called symmetric tensors of order \(n\) ; they form a sub- \(\mathbf{A}\) -module of \(\mathbf{T}^{n}(\mathbf{M})\) denoted by \(\mathbf{T}\mathbf{S}^{n}(\mathbf{M})\) ; we have \(\mathbf{T}\mathbf{S}^{0}(\mathbf{M})=\mathbf{A}\) , \(\mathbf{T}\mathbf{S}^{1}(\mathbf{M})=\mathbf{M}\) . We shall put \(\mathbf{T}\mathbf{S}(\mathbf{M}) = \bigoplus_{n=0}^{\infty} \mathbf{T}\mathbf{S}^{n}(\mathbf{M})\) ; this is a graded sub-A-module of \(\mathbf{T}(\mathbf{M})\) . For every

\(z \in \mathbf{T}^n(\mathbf{M})\) the element. \(\sum_{\sigma \in \mathfrak{S}_n} \sigma \cdot z\) belongs to \(\mathbf{T}\mathbf{S}^n(\mathbf{M})\) ; we denote it by

\begin{enumerate}
\def\labelenumi{\alph{enumi}.}
\setcounter{enumi}{18}
\tightlist
\item
  z and call it the symmetrization of z. The mapping \(s: z \mapsto s\) . z is a homomorphism of the A-module \(\mathbf{T}^{n}(\mathbf{M})\) into the A-module \(\mathbf{T}\mathbf{S}^{n}(\mathbf{M})\) . If \(z \in \mathbf{T}\mathbf{S}^{n}(\mathbf{M})\) , then \(s \cdot z = n! z\) .
\end{enumerate}

